\section{Introduction}
\label{sec:intro}

Automated market makers let anyone earn trading fees without quoting prices by
hand. Uniswap~v3 sharpened the offer: a liquidity provider (LP) chooses the
price range it is willing to trade in, so a deposit of \usdc{100000} earns the
fee share of a larger full-range position while the price stays inside the
chosen band\cite{adams2021uniswap}. Narrowing the band raises the fee earned
per dollar of capital, and that promise drew liquidity into v3.

The promise hides a cost. Liquidity quoted away from the market price is an
offer that arbitrageurs accept, and every such trade pushes the pool toward the
market price while leaving the LP holding more of the asset that is falling.
Because the loss is a transfer rather than a fluctuation, it does not reverse
when the price returns. Milionis et al.\ named this adverse selection cost
loss-versus-rebalancing (LVR) and showed that a full-range position loses value
at a rate proportional to the square of volatility\cite{milionis2022lvr}.
Fritsch and Canidio measured realised arbitrage losses on large pools and found
that they exceed the fees those pools earn\cite{fritsch2024measuring}, and
Loesch et al.\ showed that concentration magnifies the loss once the price
leaves the band\cite{loesch2021impermanent}. Fee income alone is not the LP's
return.

An LP who wants more than a passive range has three decisions, and they
interact. When to recenter trades fee income against gas and slippage; how wide
to set the next band trades fee density against the time until an edge; how
much of the position to hedge trades directional risk against funding and perp
fees. Each decision is an impulse, costly and taken under uncertainty about
volatility, order flow and regime. The objective is not to earn the most but to
earn the most per unit of risk taken, which places the problem within
risk-sensitive impulse control rather than static portfolio
choice\cite{anchuri2026rammstein}.

Recent work takes up parts of it. Zhang et al.\ train a hedged agent for v3
provision, but the hedge executes without cost or funding and the policy sees
no regime information\cite{zhang2023adaptive}. ZeroSwap learns to set a market
maker's spread so that quoted prices track an external reference, a different
lever from range placement\cite{nadkarni2023zeroswap}. RAmmStein derives
regime-dependent rebalancing thresholds and reports fewer rebalances for the
same net return, but its net metric excludes impermanent loss and it works with
off-chain data for a single pair\cite{anchuri2026rammstein}. Chionas et al.\
study risk aversion in concentrated liquidity inside a simulation and leave out
on-chain data and the hedge\cite{chionas2026concentrated}. The closest study
learns when to redeploy and how wide to stay on the same primary pool used
here, but it leaves the hedge out and reports cumulative reward
only\cite{xu2025defi}. What none of these provides is the combination this
report is designed to test: one risk-sensitive policy over timing, width and
hedge, every friction charged on both legs, evaluated against hand-written
rules on the trades that actually occurred.

This report fills that gap. We build a testbed on three years of recorded
trades from two real ETH/USDC pools, with an exact replay backtester that
computes fees in the protocol's own integer arithmetic, a fast simulator for
training, and a parity harness that keeps the two honest on identical inputs.
On that testbed we ask four questions.

\begin{itemize}
\item \textbf{Where do fees fail?} We decompose passive returns into fees,
impermanent loss, gas and slippage, and locate the conditions under which the
losses win.
\item \textbf{Why does naive learning fail?} We document the collapse of a
standard policy on this task, trace it to three interface problems, and measure
what each repair buys.
\item \textbf{Can one policy do better?} We train a risk-sensitive policy over
timing, width and hedging and compare it out of sample with the best
hand-written rules and a hedged benchmark.
\item \textbf{Where is active management worth it?} We map the boundary between
active and passive provision across volatility, fee tier, gas and capital.
\end{itemize}

The failures of the first policies are part of the result, not a footnote. In
this domain, what the agent is allowed to see and what it is rewarded for
matter more than the choice of learning algorithm, and a report that showed
only the eventual policy would hide the more transferable lesson.
